% ============================================================
%  [*] About this document — what it covers and how to read it
% ============================================================
\chapter*{About this document}
\addcontentsline{toc}{chapter}{About this document}
\markboth{About this document}{About this document}

This document describes the \product{} (\lt{\productlt}) as it was deployed on \docdate: what
each part does, why it is built that way, and how the parts fit together. It is written for a
reader who knows compliance work and financial regulation and wants to understand the system
--- above all how it uses artificial intelligence --- without reading the code. It names the
files and functions as it goes, so that a reader who does open the code finds the place.

\section*{How it is arranged}

\begin{description}[style=nextline]
  \item[Part I --- The system at a glance] What the system is for, the services it runs as,
    and the data it keeps (\cref{ch:intro,ch:arch,ch:data}).
  \item[Part II --- Artificial intelligence in the system] The longest part. It follows the
    data from the law to the report: how a model is reached at all and what a call costs
    (\cref{ch:models}), how prompts are kept and changed (\cref{ch:prompts}), how the legal
    texts become a searchable corpus (\cref{ch:corpus}) and a register of obligations
    (\cref{ch:obligations}), what is read out of a company's documents (\cref{ch:company}),
    how the compliance report is made (\cref{ch:report}) and lives on (\cref{ch:register}),
    how the documents of a compliance programme are drafted (\cref{ch:programme}), and how the
    chat assistant works (\cref{ch:chat}). \Cref{ch:control} then gathers what keeps all of
    this trustworthy: the checks, the evaluation, the people who decide, and the cost.
  \item[Part III --- Using and running the system] The console page by page
    (\cref{ch:console}), the API (\cref{ch:api}), security (\cref{ch:security}), deployment
    and operations (\cref{ch:ops}), the tests (\cref{ch:testing}), and how the system answers
    the course it was built for (\cref{ch:course}).
  \item[Reference] Every prompt (\cref{app:prompts}), every chat tool (\cref{app:tools}),
    every act of the catalogue (\cref{app:acts}), the settings (\cref{app:config}) and a
    glossary (\cref{app:glossary}).
\end{description}

\section*{Who is acting}

A system that uses language models is only as trustworthy as the line it draws between what a
model wrote, what the code verified and what a person decided. The document marks that line
the same way everywhere.

\begin{aibox}{A model is asked}
A yellow panel describes a call to a language model: what it is given, what it is told, what it
must answer. In the figures a yellow box is a model call.
\end{aibox}

\begin{codebox}{Code checks or computes}
A blue panel describes something done without a model: a quotation located in its document, a
provision looked up in the corpus, a score computed, a rule enforced. Nothing in such a panel
depends on what a model happened to say.
\end{codebox}

\begin{humanbox}{A person decides}
A framed panel marks a decision the system never takes itself: approving the scope of a
report, approving an obligation, overriding a finding, signing a report off, approving a
drafted policy.
\end{humanbox}

\begin{lessonbox}{Measured}
A grey panel reports a measurement or a failure, with its date and what was changed because of
it. The figures in these panels are observations made on the course's sample company on the
day named; they are not promises about other companies or other days.
\end{lessonbox}

Every place where a model is called also has a \emph{path card}: one table with the same ten
rows --- the trigger, the input, the prompt and the model, the output, what code checks, what is
stored, what happens on failure, the cost, the human control, and how the path is measured.
The cards are listed in \cref{ch:control}.

\section*{Conventions}

\begin{tabularx}{\textwidth}{@{}L{44mm}Y@{}}
  \toprule
  \hd{Written as} & \hd{Means} \\
  \midrule
  \file{analysis/pipeline.py} & A file. Paths of the API are given from \file{django/compliance/}, of the assistant service from \file{assistant/app/}, of the console from \file{vite/app/src/}; any other path is given from the root of the repository. \\
  \code{run_prompt} & A name in the code: a function, a constant, a field, a value. \\
  \prompt{report.gap_analysis} & A prompt, by the name it is stored under. \\
  \model{gpt-6.1-sol} & A model, by the name the gateway gives it. \\
  \route{POST}{/api/reports} & A route of the API. \\
  \ui{Veiksmai} & A label as the console prints it. \\
  \lt{atitikties pareigūnas} & Lithuanian words. The product's first language is Lithuanian; it also speaks English. \\
  \status{met} \status{partial} \status{missing} \status{unknown} \status{na} & The five statuses of a finding. \\
  \bottomrule
\end{tabularx}

Amounts are in US dollars, as the model gateway bills them. Times are UTC.
